\documentclass[10pt,a4paper]{amsart}
\usepackage[margin=19mm]{geometry}
\usepackage{amsmath,amssymb,mathtools}
\usepackage[scr=rsfs,cal=euler]{mathalfa}
\usepackage{xfrac}
\usepackage[unicode,hidelinks,bookmarks=false]{hyperref}
\newcommand{\ie}{i.e.\ }
\newcommand{\cf}{cf.\ }
\newcommand{\resp}{respectively\ }
\newcommand{\Subsection}{\subsection}
\providecommand{\nobreakdash}{\nobreak\hbox{-}}
\setlength{\emergencystretch}{2em}
\multlinegap0pt
\usepackage{bm}
\usepackage{bbm}
\usepackage{dsfont}
\usepackage{xspace}
\usepackage{cancel}
\usepackage{mathtools}
\usepackage{amsthm}
\makeatletter
\@ifundefined{case}{\newtheorem{case}{Case}}{}
\@ifundefined{claim}{\newtheorem{claim}{Claim}}{}
\@ifundefined{corollary}{\newtheorem{corollary}{Corollary}}{}
\@ifundefined{lemma}{\newtheorem{lemma}{Lemma}}{}
\@ifundefined{proposition}{\newtheorem{proposition}{Proposition}}{}
\makeatother
\newcommand{\MMM}{\mathbf M}
\newcommand{\RR}{\mathbf{R}}  % reals
\newcommand{\Tan}{\operatorname{Tan}}
\newcommand{\eps}{\epsilon}
\newcommand{\inj}{\operatorname{inj}}    %injectivity radius
\title[The Avoidance Principle for Noncompact Hypersurfaces Moving ]{The Avoidance Principle for Noncompact Hypersurfaces Moving by Mean Curvature Flow}
\author{Brian White}
\date{Source: arXiv:2401.13966v3 (CC BY 4.0)}
\begin{document}
\maketitle

\noindent\emph{Source and licence.} The Avoidance Principle for Noncompact Hypersurfaces Moving by Mean Curvature Flow. Version arXiv:2401.13966v3 (CC BY 4.0), distributed under the Creative Commons Attribution 4.0 International licence, as recorded in the arXiv licence metadata for this exact version and shown on the abstract page https://arxiv.org/abs/2401.13966v3. Full rights notice: licenses/arxiv-2401.13966-CC-BY-4.0.txt.\par\smallskip

\noindent\textit{Editorial scope.} Counted unit: the Claim whose source label is \texttt{claim-two} in the article (source0.tex lines 668--678, source label \texttt{claim-two}), delivered together with the article's own complete original proof of it (source0.tex lines 680--803, one \texttt{proof} environment of 124 lines). No mathematics is written, repaired or re-derived: statement and proof are the article's own text line for line. Required context retained uncounted: speed-lemma (lines 336--346); double-speed-corollary (lines 364--374); key-fact (lines 442--456); case-1 (lines 533--534); claim-one (lines 583--589). Every definition, display and label of those lines is retained unchanged. Omitted from this package and not redistributed: 1--335, 347--363, 375--441, 457--532, 535--582, 590--667, 679--679, 804--1865. Nothing inside a retained excerpt is omitted or altered: every retained body is delivered exactly as the article's lines are written, so no omission receipt of any kind accompanies this package. Citations: the retained excerpts carry no citation command at all, so no bibliography entry of the article is transcribed and no reference is redistributed. Reference adaptation: none was needed and none was made; every reference inside the delivered unit resolves inside it, so no unresolved reference or citation is produced. Printed slips kept literal: no printed word, symbol, hypothesis or number of the counted statement or of its proof was altered. Layout and class: the article's own class is \texttt{amsart} with packages including chngcntr, apptools, color, amsthm, amssymb, verbatim, graphicx, enumerate, hyperref, amsrefs; the wrapper is the lane's pinned \texttt{amsart} editorial wrapper at 10pt on A4, which restates the theorem-like environments the retained text needs and adds the article's own macro definitions that the retained text needs (\texttt{\textbackslash MMM}, \texttt{\textbackslash RR}, \texttt{\textbackslash Tan}, \texttt{\textbackslash eps}, \texttt{\textbackslash inj}), with extra wrapper packages (bm, bbm, dsfont, xspace, cancel, mathtools). The delivered numbering is the unit's own. Assets: no retained span contains a figure environment, a graphics-inclusion command, a PDF-page inclusion, a table or a TikZ picture; no image file is copied into this package, the package declares no asset dependency and no image file is redistributed with it. Licence: version arXiv:2401.13966v3 of this article is distributed under the Creative Commons Attribution 4.0 International licence, as recorded in the arXiv licence metadata for this exact version and shown on the abstract page https://arxiv.org/abs/2401.13966v3. A full rights notice is delivered at licenses/arxiv-2401.13966-CC-BY-4.0.txt. Characters: every retained excerpt is pure ASCII, so the delivered engine, XeLaTeX with its default Latin Modern text font, reports no missing character.
\begin{lemma}[Finite Speed Lemma]\label{speed-lemma}
For $r>0$, there is an $h=h(r,\Lambda,m)$ with the following property.
If $S\subset N$, if $t\in [0,\infty)\mapsto X(t)$ is a weak set flow, and if
\[
    r < R:=d(X(0),S),
\]
then
\[
    d(X(t),S) \ge R - ht  \qquad\text{for $0\le t \le (R-r)/h$}.
\]
\end{lemma}

\begin{corollary}\label{double-speed-corollary}
If $t\in [t_0,\infty)\mapsto X(t), \, Y(t)$ are weak set flows with
\[
   d(X(t_0),Y(t_0))  > r > 0,
\]
then
\[
    d(X(t_0+t),Y(t_0+t)) \ge d(X(t_0), Y(t_0))  - 2 h t
\]
for $t\le (R-r)/h$, where $h$ is as in Lemma~\ref{speed-lemma}.
\end{corollary}

\begin{proposition}\label{key-fact}
Suppose that $\lambda\in \RR$,  
 that $W$ is a smooth (not necessarily complete)
Riemannian manifold, and 
    that $X(\cdot)$ and $Y(\cdot)$ are weak set flows in $W$. 
  Suppose that $T>0$ and that
\begin{align*}
   e^{-\lambda t} d(X(t),Y(t)) &> \eta \quad\text{for $0\le t<T$}, \\
   e^{-\lambda T} d(X(T),Y(T)) &= \eta.
\end{align*}
Suppose also that there is a geodesic $\Gamma$ of length $d(X(T),Y(T))$
from a point $x\in X(T)$ to a point $y\in Y(T)$, 
and that $(y,T)$ is a regular point of the flow $Y(\cdot)$.
Then there are points in $W$  (indeed, points on $\Gamma$) where the Ricci curvature is $\le \lambda$.
\end{proposition}

\begin{case}\label{case-1} $d(X(0),Y(0)) \le  r:= \frac12 \inj(N)$.
\end{case}

\begin{claim}\label{claim-one}
There is an $\eps>0$ with the following property.
If $0<T\le \eps$, if $p\in \MMM(T)$, and if $d(p,K^c)\le T^{1/2}$, then
\[
  T^{1/2}   A(\MMM,p,T) < 1.
\]
\end{claim}

\begin{claim}\label{claim-two}
There is a $\delta>0$ such that 
\[
    e^{-\lambda t} d(X(t),\MMM(t)) \ge R
\]
and
\[
   e^{-\lambda t} d(Y(t),\MMM(t)) \ge R
\]
for $0\le t\le \delta$.
\end{claim}

\begin{proof}
It suffices to prove it for $X(\cdot)$.
Recall that there exist $C<\infty$ and $\eta<0$ such that
\begin{equation}\label{C-eta}
    \text{$d(X(t), K) \ge R - Ct$ for $0\le t\le \eta$.}
\end{equation}
(See Lemma~\ref{speed-lemma}.)
Choose $\delta>0$ small enough that
\begin{align*}
\delta&\le \eps, \\
\delta&\le \eta, \\
C\delta^{1/2} &\le \frac12
\end{align*}
where $\eps$ is as in Claim~\ref{claim-one} and $\eta$ and $C$ are as in~\eqref{C-eta}.
Thus if $0\le t \le \delta$, then
\begin{equation}\label{Ct-small}
  Ct \le C \delta^{1/2} t^{1/2}  \le  \frac12 t^{1/2}.
\end{equation}
Let 
\begin{align*}
&f: [0,\delta] \to \RR, \\
&f(T) = \inf_{0\le t\le T} e^{-\lambda t} d(X(t),\MMM(t)).
\end{align*}
Suppose the claim is false.  Then $f(T)<R$ for some $T\in (0,\delta]$.
 By Corollary~\ref{double-speed-corollary}, $f$
is continuous.
Thus, by replacing $T$ by a smaller $T>0$, we can
assume that 
\[
   0 < f(T) < R.
\]
We can also assume that $f(t)>f(T)$ for $t<T$; otherwise replace $T$
by the smallest $t$ such that $f(t)=f(T)$.
Thus
\[
   \eta:=f(T) < f(t) \qquad\text{for $t<T$}.
\]
In particular, $e^{-\lambda T} d(X(T),\MMM(T)) =\eta$,
so
\begin{equation}\label{close-enough}
d(X(T),\MMM(T)) = e^{\lambda T}\eta < \eta < R \le \frac12r
\end{equation}
since we are assuming that $\lambda<0$ and (in Case~\ref{case-1}) that $R\le \frac12r$.




Choose $p_i\in \MMM(T)$ and $q_i\in X(T)$ so that
\[
   d(p_i,q_i) \to d(\MMM(T),X(T)) = e^{\lambda T}\eta.
\]

By passing to a subsequence, we can assume that $d(K^c,q_i)$ converges to a limit in $[0,\infty]$.
We wish to show that
\begin{equation}\label{the-inequality}
  \lim_i d(q_i, K^c) \le \frac12 T^{1/2}.
\end{equation}
We may assume that $q_i$ is in the interior of $K$ for all sufficiently large $i$, as otherwise
the inequality~\eqref{the-inequality} is trivially true. 
For such $i$,
\begin{align*}
d(p_i,q_i)
&\ge 
d(p_i, \partial K) + d(\partial K,q_i) \\
&\ge
R- CT + d(K^c,q_i).
\end{align*}
Thus
\[
e^{\lambda T}\eta \ge R - CT + \lim_i d(K^c,q_i).
\]
Now $e^{\lambda T} \le 1$ (since we are assuming that $\Lambda\le 0$)
 and $\eta<R$, so
\[
   \lim_i d(q_i,K^c) \le CT < \frac12 T^{1/2}
\]
by~\eqref{Ct-small}.
This completes the proof that the inequality~\eqref{the-inequality} holds.

By  Claim~\ref{claim-one} and~\eqref{the-inequality},
\begin{equation}\label{smooth-point}
 |T|^{1/2} A(\MMM_i,q_i,T) < 1
\end{equation} 
for all sufficiently large $i$.


Now let $L_i: \RR^{m+1} \to \Tan(N,p_i)$ be a a linear isometry,
and let
\begin{align*}
&F_i: B=B^{m+1}(0,r) \to N, \\
&F_i = \exp_{p_i}\circ L_i.
\end{align*}
Let $g_i$ be the pull-back by $F_i$ of the metric on $N$, so that
$F_i$ maps $(B,g_i)$ isometrically onto $B(p_i,r)$.

For large $i$, $d(p_i,q_i) < r$ (by~\eqref{close-enough}),
so there is a $\tilde q_i\in B$ with $F_i(\tilde q_i) = q_i$.
Let 
\begin{align*}
\tilde X_i(t) &= F_i^{-1}(X(t)), \\
\tilde \MMM_i(t) &= F_i^{-1}(\MMM(t)).
\end{align*}
Thus $\tilde X_i(\cdot)$ and $\tilde \MMM_i(\cdot)$ are weak set flows in $B$ with respect to $g_i$.

After passing to a subsequence, the $g_i$ converge to a smooth Riemannian metric $\tilde g$,
  $\tilde X_i(\cdot)$ and $\tilde \MMM_i(\cdot)$ 
converge to weak set flows $\tilde X(\cdot)$ and $\tilde \MMM(\cdot)$ in $(B,\tilde g)$,
 and $\tilde q_i$ converges to point $\tilde q$ with
 \begin{equation}\label{attained}
    d(0,\tilde q) = |\tilde q| = d(X(T),\MMM(T))= e^{\lambda T}\eta.
\end{equation}

For $0\le t<T$, 
\begin{equation}\label{not-attained}
  e^{-\lambda t} d(\tilde X(t), \tilde \MMM(t)) \ge  e^{-\lambda t}d(X(t), \MMM(t)) > \eta.
\end{equation}

By~\eqref{smooth-point},
 the spacetime point $(0,T)$ is a regular point of the flow $\tilde \MMM$.
Thus by~\eqref{attained}, \eqref{not-attained}, and~Proposition~\ref{key-fact},
    $(B,\tilde g)$ has points of Ricci curvature $\le\lambda$,
contrary to our choice of $\lambda<\Lambda$.  
This completes the proof of Claim~\ref{claim-two}.
\end{proof}

\end{document}
